% from ega1-1-fr.tex
\oldpage[I]{96}
où $j:N'\to N'\otimes_{A'}A_{[\varphi]}$ est l'homomorphisme
$z'\mapsto z'\otimes 1$. De même, l'application canonique
(\hyperref[0.4.4.3.3-fr]{0, 4.4.3.3})
$\sigma:\Phi^*(\Phi_*(\widetilde{M}))\to\widetilde{M}$ n'est autre que
$\widetilde{p}$, où
$p:M_{[\varphi]}\otimes_{A'}A_{[\varphi]}\to M$ est l'homomorphisme
canonique qui, à tout produit tensoriel $z\otimes a$ ($z\in M$, $a\in A$)
fait correspondre $a\cdot z$~; cela résulte aussitôt des définitions
(\hyperref[0.3.7.1-fr]{0, 3.7.1},
\hyperref[0.4.4.3-fr]{0, 4.4.3}, et
\hyperref[I.1.3.7-fr]{I, 1.3.7}).

On en conclut (\hyperref[0.4.4.3-fr]{0, 4.4.3} et
\hyperref[0.3.5.4.4-fr]{3.5.4.4}) que si
$v:N'\to M_{[\varphi]}$ est un $A'$-homomorphisme, on a
$\widetilde{v}^{\#}=(v\otimes 1)^{\sim}$.
\end{env}

\begin{env}[1.6.8]
\label{I.1.6.8-fr}
Soient $N'_1$, $N'_2$ deux $A'$-modules, $N'_1$ étant supposé avoir une
\emph{présentation finie}~; il résulte alors de
(\hyperref[I.1.6.7-fr]{1.6.7}) et
(\hyperref[I.1.3.12-fr]{1.3.12, (ii)}) que l'homomorphisme canonique
(\hyperref[0.4.4.6-fr]{0, 4.4.6})
\[
  \Phi^*\bigl(\mathscr{H}\!om_{\widetilde{A'}}
    (\widetilde{N'_1},\widetilde{N'_2})\bigr)
  \longrightarrow
  \mathscr{H}\!om_{\widetilde A}
    \bigl(\Phi^*(\widetilde{N'_1}),\Phi^*(\widetilde{N'_2})\bigr)
\]
n'est autre que $\widetilde{\gamma}$, en désignant par $\gamma$
l'homomorphisme canonique de $A$-modules
\[
  \operatorname{Hom}_{A'}(N'_1,N'_2)\otimes_{A'}A
  \longrightarrow
  \operatorname{Hom}_A(N'_1\otimes_{A'}A,N'_2\otimes_{A'}A).
\]
\end{env}

\begin{env}[1.6.9]
\label{I.1.6.9-fr}
Soient $\mathfrak{J}'$ un idéal de $A'$, $M$ un $A$-module~; comme par
définition $\widetilde{\mathfrak{J}'}\widetilde M$ est l'image de
l'homomorphisme canonique
$\Phi^*(\widetilde{\mathfrak{J}'})\otimes_{\widetilde A}\widetilde M
\to\widetilde M$, il résulte de (\hyperref[I.1.6.5-fr]{1.6.5}) et
(\hyperref[I.1.3.12-fr]{1.3.12, (i)}) que
$\widetilde{\mathfrak{J}'}\widetilde M$ s'identifie canoniquement à
$(\mathfrak{J}'M)^{\sim}$~; en particulier,
$\Phi^*(\widetilde{\mathfrak{J}'})\widetilde A$ s'identifie à
$(\mathfrak{J}'A)^{\sim}$, et en tenant compte de l'exactitude à droite du
foncteur $\Phi^*$, la $\widetilde A$-Algèbre
$\Phi^*((A'/\mathfrak{J}')^{\sim})$ s'identifie à
$(A/\mathfrak{J}'A)^{\sim}$.
\end{env}

\begin{env}[1.6.10]
\label{I.1.6.10-fr}
Soient $A''$ un troisième anneau, $\varphi'$ un homomorphisme
$A''\to A'$, et posons $\varphi''=\varphi\circ\varphi'$. Il résulte
aussitôt des définitions que
${}^a\varphi''=({}^a\varphi')\circ({}^a\varphi)$, et
$\widetilde{\varphi''}=\widetilde\varphi\circ\widetilde{\varphi'}$
(\hyperref[0.1.5.7-fr]{0, 1.5.7}). On en conclut que l'on a
$\Phi''=\Phi'\circ\Phi$~; en d'autres termes,
$(\operatorname{Spec} A,\widetilde A)$ est un \emph{foncteur} en $A$ de la
catégorie des anneaux dans celle des espaces annelés.
\end{env}

\subsection{Caractérisation des morphismes de schémas affines.}
\label{subsection:I.1.7-fr}

\begin{definition}[1.7.1]
\label{I.1.7.1-fr}
On dit qu'un espace annelé $(X,\mathscr{O}_X)$ est un \emph{schéma affine}
s'il est isomorphe à un espace annelé de la forme
$(\operatorname{Spec}(A),\widetilde A)$, où $A$ est un anneau~; on dit alors
que $\Gamma(X,\mathscr{O}_X)$, qui s'identifie canoniquement à l'anneau $A$
(\hyperref[I.1.3.7-fr]{1.3.7}) est l'anneau du schéma affine
$(X,\mathscr{O}_X)$, et on le note $A(X)$ quand aucune confusion n'en résulte.
\end{definition}

Par abus de langage, quand nous parlerons du \emph{schéma affine}
$\operatorname{Spec}(A)$, il s'agira toujours de l'espace annelé
$(\operatorname{Spec}(A),\widetilde A)$.

\begin{env}[1.7.2]
\label{I.1.7.2-fr}
Soient $A$, $B$ deux anneaux, $(X,\mathscr{O}_X)$,
$(Y,\mathscr{O}_Y)$ les schémas affines correspondants sur les spectres
premiers $X=\operatorname{Spec}(A)$, $Y=\operatorname{Spec}(B)$. On a vu
(\hyperref[I.1.6.1-fr]{1.6.1}) qu'à tout homomorphisme d'anneaux
$\varphi:B\to A$ correspond un morphisme
$\Phi=({}^a\varphi,\widetilde\varphi)=\operatorname{Spec}(\varphi):
(X,\mathscr{O}_X)\to(Y,\mathscr{O}_Y)$. On notera que $\varphi$ est
entièrement déterminé par $\Phi$, car on a par définition
\[
  \varphi=\Gamma(\widetilde\varphi):
  \Gamma(\widetilde B)\longrightarrow
  \Gamma({}^a\varphi_*(\widetilde A))=\Gamma(\widetilde A).
\]
\end{env}

\begin{theorem}[1.7.3]
\label{I.1.7.3-fr}
Soient $(X,\mathscr{O}_X)$, $(Y,\mathscr{O}_Y)$ deux schémas affines. Pour
qu'un morphisme d'espaces annelés
$(\psi,\theta):(X,\mathscr{O}_X)\to(Y,\mathscr{O}_Y)$ soit de la forme
$({}^a\varphi,\widetilde\varphi)$, où $\varphi$ est un homomorphisme
d'anneaux $A(Y)\to A(X)$, il faut et il suffit que, pour tout $x\in X$,
$\theta_x^{\#}$ soit un homomorphisme local~:
$\mathscr{O}_{\psi(x)}\to\mathscr{O}_x$.
\end{theorem}

